\documentclass[11pt, a4paper]{article}

% ── Packages ──────────────────────────────────────────────────────────
\usepackage[utf8]{inputenc}
\usepackage[T1]{fontenc}
\usepackage{lmodern}
\usepackage[margin=1in]{geometry}
\usepackage{amsmath, amssymb, amsthm}
\usepackage{booktabs}
\usepackage{longtable}
\usepackage{graphicx}
\usepackage{hyperref}
\usepackage{xcolor}
\usepackage{listings}
\usepackage{tcolorbox}
\usepackage{polya}

\hypersetup{
  colorlinks=true,
  linkcolor=polyablue,
  urlcolor=polyablue,
  citecolor=polyablue,
}

% ── Code listing style ────────────────────────────────────────────────
\lstdefinestyle{polya-python}{
  language=Python,
  basicstyle=\ttfamily\small,
  keywordstyle=\color{polyablue}\bfseries,
  commentstyle=\color{polyagray}\itshape,
  stringstyle=\color{polyagreen},
  numbers=left,
  numberstyle=\tiny\color{polyagray},
  numbersep=8pt,
  frame=single,
  framerule=0.4pt,
  rulecolor=\color{polyagray},
  breaklines=true,
  showstringspaces=false,
  tabsize=4,
}
\lstset{style=polya-python}
\title{Customer Churn Classification}
\author{uber-polya}
\date{2026-02-26}

\begin{document}

\maketitle
\tableofcontents
\newpage

% ── Phase A: Problem Formulation ─────────────────────────────────────
\section{Problem Formulation}

\begin{polyamodel}

\textbf{Problem Type}: Find \hfill
\textbf{Domain}: Machine Learning


\end{polyamodel}

% ── Universe ──────────────────────────────────────────────────────────
\subsection{Universe}

\begin{itemize}
  \item $Best Model: \{name, accuracy, roc_auc\}$
  \item $Confusion Matrix: \{TP, FP, FN, TN\}$
  \item $Feature Importances: \{tenure, total_charges, monthly_charges, num_support_tickets, contract_type, has_internet\}$
  \item $Cross Validation: \{n_folds, scores, mean, std\}$
\end{itemize}

% ── Variables ─────────────────────────────────────────────────────────
\subsection{Variables}

\begin{longtable}{lll}
\toprule
\textbf{Symbol} & \textbf{Meaning} & \textbf{Type / Range} \\
\midrule
\endhead
$x$ & decision variable & $\mathbb{{R}}$ \\
\bottomrule
\end{longtable}

% ── Structure ─────────────────────────────────────────────────────────
\subsection{Structure}

- **Data**: \textasciitilde{}1,000 synthetic customers with 6 features (seed=42)
- **Features**: tenure (months), monthly charges, total charges, support tickets, contract type (month-to-month vs annual), internet service (yes/no)
- **Target**: churned (1) vs retained (0), \textasciitilde{}30\% churn rate
- **Churn drivers**: shorter tenure, more support tickets, and month-to-month contracts increase churn probability
- **Question**: Which model best predicts churn, and which features matter most?

% ── Mapping ───────────────────────────────────────────────────────────
\subsection{Real-World Mapping}

\begin{longtable}{ll}
\toprule
\textbf{Real-World Concept} & \textbf{Mathematical Object} \\
\midrule
\endhead
problem input & $instance parameters$ \\
\bottomrule
\end{longtable}

% ── Constraints ───────────────────────────────────────────────────────
\subsection{Constraints}

\begin{enumerate}
  \item $Accuracy**: fraction of correct predictions (TP + TN) / total$
  \hfill \textit{()}
  \item $Precision**: of those predicted as churned, how many actually churned (TP / (TP + FP))$
  \hfill \textit{()}
  \item $Recall**: of those who actually churned, how many were caught (TP / (TP + FN))$
  \hfill \textit{()}
  \item $F1 score**: harmonic mean of precision and recall$
  \hfill \textit{()}
  \item $ROC-AUC**: area under the receiver operating characteristic curve; measures discrimination across all thresholds$
  \hfill \textit{()}
  \item $Confusion matrix**: 2x2 table of TP, FP, FN, TN$
  \hfill \textit{()}
  \item $Feature importance**: mean decrease in impurity (Gini) across all trees in the Random Forest$
  \hfill \textit{()}
  \item $Cross-validation**: 5-fold stratified CV estimates out-of-sample performance and its variability$
  \hfill \textit{()}
\end{enumerate}

% ── Objective / Claim ─────────────────────────────────────────────────
\subsection{Claim}


% ── Approach ──────────────────────────────────────────────────────────
\subsection{Approach}

\begin{description}
  \item[Suggested approach:] Logistic Regression + Random Forest (100) + Gradient Boosting (100)
  \item[Complexity class:] 
  \item[Available tools:] Logistic Regression + Random Forest
\end{description}
% ── Phase B: Solution ────────────────────────────────────────────────
\section{Solution}

\begin{polyaresult}

\textbf{Answer}: Solution computed


\textbf{Optimal}: No / Unknown \hfill
\textbf{Feasible}: No
\textbf{Algorithm}: Logistic Regression + Random Forest (100) + Gradient Boosting (100) \quad $$

\textbf{Time}: 0.34118065395159647s


\end{polyaresult}

% ── Solution Details ──────────────────────────────────────────────────
\subsection{Solution Details}


Method: 1. Logistic Regression: Linear model with L2 regularization. Fast baseline. Provides interpretable coefficients for each feature.

2. Random Forest (100 trees): Bagged ensemble of decision trees. Handles non-linear relationships and feature interactions. Produces feature importance rankings.

3. Gradient Boosting (100 trees): Sequential boosting of shallow trees. Often achieves the best accuracy on tabular data. Minimizes log-loss via gradient descent in function space.

% ── Verification ──────────────────────────────────────────────────────
\subsection{Verification}

\begin{longtable}{lcl}
\toprule
\textbf{Check} & \textbf{Status} & \textbf{Detail} \\
\midrule
\endhead
Best Accuracy Gt 75Pct & \passmark & Passed \\
Best Roc Auc Gt 70Pct & \passmark & Passed \\
Rf Outperforms Baseline & \passmark & Passed \\
Feature Importances Sum To 1 & \passmark & Passed \\
Cv Accuracy Stable & \passmark & Passed \\
Confusion Matrix Sums Correct & \passmark & Passed \\
All Passed & \passmark & Passed \\
\bottomrule
\end{longtable}

% ── Phase C: Interpretation ──────────────────────────────────────────
\section{Interpretation}

% ── The Question & The Answer ─────────────────────────────────────────
\subsection{The Question}
- **Data**: \textasciitilde{}1,000 synthetic customers with 6 features (seed=42)
- **Features**: tenure (months), monthly charges, total charges, support tickets, contract type (month-to-month vs annual), internet service (yes/no)
- **Target**: churned (1) vs retained (0), \textasciitilde{}30\% churn rate
- **Churn drivers**: shorter tenure, more support tickets, and month-to-month contracts increase churn probability
- **Question**: Which model best predicts churn, and which features matter most?.

\subsection{The Answer}

\begin{polyainsight}[title={Bottom Line}]
A feasible solution was found.
\end{polyainsight}

% ── What This Means ───────────────────────────────────────────────────
\subsection{What This Means}

- Best model: Logistic Regression (accuracy \textasciitilde{}80\%, ROC-AUC \textasciitilde{}0.85)
- All models beat majority-class baseline (\textasciitilde{}70\% accuracy by always predicting "not churned")
- Top features (Random Forest importance): tenure, total charges, monthly charges, support tickets
- Cross-validation: stable accuracy (std < 0.05 across 5 folds)
- All 7 verification checks pass

1. Logistic Regression: Linear model with L2 regularization. Fast baseline. Provides interpretable coefficients for each feature.

2. Random Forest (100 trees): Bagged ensemble of decision trees. Handles non-linear relationships and feature interactions. Produces feature importance rankings.

3. Gradient Boosting (100 trees): Sequential boosting of shallow trees. Often achieves the best accuracy on tabular data. Minimizes log-loss via gradient descent in function space.

% ── Sensitivity Analysis ──────────────────────────────────────────────

% ── Figures ───────────────────────────────────────────────────────────

% ── Recommendations ───────────────────────────────────────────────────
\subsection{Recommendations}

\begin{enumerate}
  \item Review the model assumptions to ensure real-world fidelity.
\end{enumerate}

% ── Limitations ───────────────────────────────────────────────────────
\subsection{Limitations}

\begin{polyawarnbox}
\begin{itemize}
  \item Model assumes deterministic parameters; real-world variability not captured.
  \item Solution quality depends on the accuracy of input data.
\end{itemize}
\end{polyawarnbox}

% ── Appendix ─────────────────────────────────────────────────────────

\end{document}